\section*{ID6710: Numeric Recipes \& Quick Example Calculations: F\#}
\paragraph{Metadata.}
PiID: 1913933918568 | RTEID: RTE:P26056.4201:T5.2339 | UUID: ec2b70b5-139b-5be1-ac6a-fe0e38ee4274

\paragraph{Canonical Statement.}
\# 6) Numeric recipes \& quick example calculations - **Airy radius (linear on sensor)**: \textbackslash{}(r\_\{\textbackslash{}text\{Airy\}\} = 1.22\textbackslash{},\textbackslash{}lambda\textbackslash{}, f\textbackslash{}\#\textbackslash{}). Example: \textbackslash{}(\textbackslash{}lambda=550\textbackslash{} \textbackslash{}text\{nm\}\textbackslash{}), \textbackslash{}(f\textbackslash{}\#=4\textbackslash{}) → \textbackslash{}(r\_\{\textbackslash{}text\{Airy\}\} \textbackslash{}approx 1.22\textbackslash{}times 550\textbackslash{}times10\textasciicircum{}\{-9\}\textbackslash{}times4\textbackslash{}approx 2.7\textbackslash{}times10\textasciicircum{}\{-6\}\textbackslash{}) m = 2.7 µm (compare to your pixel size). - **Angular resolution**: \textbackslash{}(\textbackslash{}theta\_1 \textbackslash{}approx 1.22\textbackslash{},\textbackslash{}lambda/D\textbackslash{}). If you estimate \textbackslash{}(\textbackslash{}theta\_1\textbackslash{}) from the image, you get \textbackslash{}(D\textbackslash{}). - **Harmonic cutoff**: \textbackslash{}(\textbackslash{}ell\_\{\textbackslash{}max\}\textbackslash{}approx 2.57\textbackslash{},D/\textbackslash{}lambda\textbackslash{}). For \textbackslash{}(D/\textbackslash{}lambda = 1000\textbackslash{}) → \textbackslash{}(\textbackslash{}ell\_\{\textbackslash{}max\}\textbackslash{}approx 2570\textbackslash{}) (very fine angular detail). - **Vane spike model**: single vane width \textbackslash{}(w\textbackslash{}) → FT amplitude \textbackslash{}(\textbackslash{}propto \textbackslash{}operatorname\{sinc\}(k w \textbackslash{}sin\textbackslash{}theta)\textbackslash{}); spike angular width \textbackslash{}(\textbackslash{}sim \textbackslash{}lambda/w\textbackslash{}) (narrow vane → narrow spike).

\paragraph{Canonical Equation.}
\[
f\#=4
\]

\paragraph{Dictionary.}
Numeric Recipes \& Quick Example Calculations: F\# — f\#=4

\paragraph{Encyclopedia.}
Numeric Recipes \& Quick Example Calculations: F\# is an algebraic definition, written f\#=4. It is an algebraic definition expressing the quantity directly in terms of the variables on its right-hand side. Within the Trawin Zero Point registry it serves as one element of the framework's mathematical narrative.
